\documentclass{beamer}
\usepackage{predavanja}


\newtheorem{izrek}{Izrek}
\newtheorem{definicija}{Definicija}



\begin{document}


% Zgornje podatke nastavite z ukazi kot v dokumentih razreda `article`.
% Več o tem, kako se naredi naslovno stran, si preberite na naslovu na naslovu:
% https://www.overleaf.com/learn/latex/Beamer
% To stran preberite do vključno razdelka "Creating a table of contents".
% Ukaz `\titlepage` deluje podobno kot ukaz `\maketitle`, ki ste ga že srečali.


\title{Matematični izrazi in uporaba paketa \texttt{beamer}}
\subtitle{Matematični izrazi in uporaba paketa \texttt{beamer}}
\institute{Fakulteta za matematiko in fiziko}
\date{}

\begin{frame}
    \titlepage
\end{frame}
% Naloga 1.3.5: pripravite kazalo vsebine.
% 1. Naslov prosojnice, s kazalom vsebine naj bo "Kratek pregled"
% 2. S pomožnim parametrom `pausesections' (v oglatih oklepajih) 
%    določite, da naj se kazalo vsebine odkriva postopoma.
%    Poglejte, kako deluje ta ukaz.
% 3. Ker ni videti preveč lepo, pomožni parameter zakomentirajte.

\begin{frame}
    \frametitle{Kratek pregled}
    \tableofcontents%[pausesections]
\end{frame}


\section{Paket \texttt{beamer}}
\input{prosojnice/1-paket-beamer.tex}

\section{Paketa \texttt{amsmath} in \texttt{amsfonts}}
\input{prosojnice/2-paketa-amsmath-amsfonts.tex}

\section[Matematika, 1. del\\\large{Analiza, logika, množice}]{Matematika, 1. del}
\input{prosojnice/3-analiza-logika-mnozice.tex}

\section{Stolpci in slike}

\section{Paket \texttt{beamer} in tabele}

\section[Matematika, 2. del\\\large{Zaporedja, algebra, grupe}]{Matematika, 2. del}

\end{document}